% Capitulo: Antes do dashboard
\chapter{Antes do dashboard}
\label{cap:origens}

Antes de existir dashboard, existia o relatório impresso. Ele chegava na mesa do
gerente alguns dias depois do fechamento do mês, em formulário contínuo, com páginas e
páginas de números alinhados em colunas. Quem precisava de uma resposta diferente
daquela que o relatório trazia fazia um pedido ao centro de processamento de dados e
esperava. Às vezes semanas.

Parece pré-história, mas boa parte dos problemas que discutimos hoje já estava ali:
números que chegam tarde, relatórios que ninguém lê, a distância entre quem tem a
pergunta e quem sabe extrair a resposta. Este capítulo volta a essa época por um motivo
prático. Quem entende de onde vêm as ideias do BI reconhece mais rápido quando uma
ferramenta nova está só dando outro nome a uma ideia velha, e quando está trazendo algo
de fato novo.

\section{Relatórios, planilhas e sistemas de apoio à decisão}
\label{sec:sad}

Nos anos 1960, as grandes empresas começaram a usar computadores de grande porte para
processar a folha de pagamento, o faturamento e o estoque. Como esses dados já estavam
na máquina, parecia natural usá-los também para informar a gestão. Surgiram os sistemas
de informação gerencial, que produziam relatórios periódicos a partir dos dados
operacionais.

A promessa era dar ao gestor toda a informação de que ele precisava. Em 1967, Russell
Ackoff publicou uma crítica que continua atual. Para ele, o problema da maioria dos
gestores não era a falta de informação relevante, e sim o excesso de informação
irrelevante \parencite{ackoff1967}. Os sistemas da época partiam do princípio de que mais
dados levariam a decisões melhores, e por isso despejavam relatórios. Ackoff argumentava
o contrário: um sistema de informação útil precisa filtrar e condensar, e para isso
precisa entender como a decisão é tomada.

Quem já viu uma empresa com duzentos dashboards e nenhuma reunião em que eles sejam
abertos reconhece o diagnóstico. A tecnologia mudou várias vezes desde então; a
tentação de produzir mais relatório em vez de melhor informação continua a mesma.

A resposta acadêmica a esse problema veio no começo dos anos 1970. \textcite{gorry1971}
separaram as decisões em estruturadas, semiestruturadas e não estruturadas, e propuseram
o nome \emph{sistemas de apoio à decisão} (SAD) para os sistemas voltados às decisões
que não se resolvem com uma regra fixa. A ideia era apoiar o julgamento do gestor, não
substituí-lo. \textcite{keen1978} desenvolveram o conceito no livro que se tornou
referência da área, e a história completa desse período está bem contada por
\textcite{power2007}.

Os SAD eram, em geral, sistemas construídos para uma decisão específica: planejar a
produção, avaliar uma carteira de investimentos, escolher rotas. Eram caros e dependiam
de especialistas. A virada que levou a análise para a mesa de qualquer pessoa veio de
outro lugar: as planilhas eletrônicas.

Quando as planilhas chegaram aos computadores pessoais, no fim dos anos 1970 e ao longo
dos anos 1980, o gestor passou a poder montar a própria análise sem pedir nada a
ninguém. Foi o primeiro \enquote{self-service} da história do BI, décadas antes de o
termo existir. E foi tão bem-sucedido que até hoje a planilha é a ferramenta de análise
mais usada nas empresas, inclusive naquelas que pagam caro por ferramentas de BI.

A planilha resolveu o problema da espera e criou outro, que vai aparecer várias vezes
neste livro: cada pessoa com a sua cópia dos dados, a sua fórmula e o seu número. O
capítulo~\ref{cap:master-data} trata de uma das consequências mais caras disso, os
cadastros que existem em versões diferentes em cada planilha.

No fim dos anos 1980, havia ainda os sistemas de informação executiva, que levavam
indicadores resumidos para a alta direção em telas simples. Eram, na prática, os avós
dos dashboards. A sopa de siglas (SIG, SAD, SIE) era grande, e foi nesse ambiente que
um nome guarda-chuva pegou.

\rascunho{
  Relato opcional: o primeiro contato com relatórios estáticos ou com uma empresa que
  rodava tudo em planilhas, e o que isso ensinou. Se não houver um caso bom, a seção se
  sustenta sem relato.
}

\section{De onde vem o termo business intelligence}
\label{sec:termo-bi}

A expressão \enquote{business intelligence} é mais antiga do que costuma se imaginar.
Ela já aparece em um livro de anedotas comerciais de 1865, para descrever um banqueiro
que conseguia vantagem sobre os concorrentes porque recebia notícias antes deles
\parencite{devens1865}. A palavra \emph{intelligence}, em inglês, tem esse sentido de
informação obtida para orientar uma ação, como em serviço de inteligência.

O uso que se aproxima do atual vem de 1958. Hans Peter Luhn, pesquisador da IBM,
publicou um artigo descrevendo um \enquote{sistema de business intelligence}: um sistema
automático que resumiria documentos, identificaria o interesse de cada área da empresa e
faria a informação certa chegar a quem precisava agir \parencite{luhn1958}. Luhn definia
inteligência como a capacidade de perceber as relações entre os fatos apresentados de
forma a orientar a ação em direção a um objetivo. \pendente{conferir a definição de Luhn
no artigo original antes de manter a paráfrase}

Vale reparar no que Luhn colocava no centro. Não era o armazenamento de dados nem o
gráfico; era a entrega da informação certa para a pessoa certa, com o objetivo de
orientar uma ação. Quase setenta anos depois, essa ainda é uma boa definição do que um
time de dados deveria fazer.

O termo só se popularizou no fim dos anos 1980. Howard Dresner, então analista do
Gartner, passou a usar \enquote{business intelligence} como guarda-chuva para os
conceitos e métodos de apoio à decisão baseados em fatos, num mercado cheio de siglas
\parencite{computerworld2006}. A autoria é disputada: o termo já aparecia em
publicações anteriores, e o próprio Dresner deu versões diferentes sobre quando e onde
começou a usá-lo. Para este livro, a disputa importa menos do que o fato de que o nome
pegou, e com ele veio uma indústria inteira.

\section{O que ficou dessa época}
\label{sec:legado-origens}

Os relatórios em formulário contínuo sumiram, os SAD viraram nota de rodapé e as siglas
foram esquecidas. Três ideias dessa época, porém, continuam valendo, e vão aparecer ao
longo do livro.

A primeira é que a pergunta de negócio vem antes do dado. Luhn falava em orientar a
ação; Gorry e Scott Morton classificavam sistemas pelo tipo de decisão que apoiavam.
Nenhum dos dois começava pela tecnologia. Projetos de dados que começam pela ferramenta,
e só depois procuram uma pergunta para responder, continuam fracassando pelo mesmo
motivo que fracassavam nos anos 1970.

A segunda é a crítica de Ackoff: mais informação não é melhor informação. Ela reaparece
quando o livro discute com que frequência os dados precisam atualizar
(capítulo~\ref{cap:frequencia}) e por que vale medir o uso dos próprios produtos de dados.
Um relatório que ninguém abre é o mesmo desperdício de 1967, só que mais barato de
produzir e, na era dos créditos e tokens, nem sempre tão barato de manter.

A terceira é que a análise sempre escapa para perto de quem decide. Aconteceu com a
planilha, aconteceu com as ferramentas de self-service (capítulo~\ref{cap:self-service})
e está acontecendo de novo com a IA (capítulo~\ref{cap:ia}). Cada vez que isso acontece,
a empresa ganha agilidade e perde consistência, e alguém precisa reconstruir a
consistência sem matar a agilidade. Esse alguém, quase sempre, é o time de dados.

O passo seguinte da história foi justamente uma tentativa de resolver a consistência de
uma vez: tirar os dados de análise de dentro dos sistemas operacionais e organizá-los num
lugar próprio. É o assunto do próximo capítulo.

\begin{leituras}
  \leitura{ackoff1967}{Um artigo curto e provocador. Os cinco pressupostos errados que
    Ackoff aponta nos sistemas de informação da época continuam aparecendo em projetos de
    dashboard.}
  \leitura{gorry1971}{O artigo que deu nome aos sistemas de apoio à decisão e que
    organizou as decisões gerenciais pelo grau de estrutura.}
  \leitura{power2007}{Um panorama acessível da história dos sistemas de apoio à decisão,
    escrito por quem acompanhou a área de perto.}
  \leitura{luhn1958}{Para ler o primeiro \enquote{sistema de business intelligence} na
    fonte e perceber como a preocupação com a pessoa certa e a ação certa já estava lá.}
\end{leituras}
